\documentclass[10.5pt]{article}
\usepackage[letterpaper, margin=0.6in, top=0.5in, bottom=0.5in]{geometry}
\usepackage{amsmath,amssymb,amsthm}
\usepackage{fancyhdr}
\usepackage{xcolor}
\usepackage{tikz}
\usepackage{enumitem}
\usepackage{tcolorbox}
\usepackage{hyperref}

% Define colors matching instructor board annotations
\definecolor{chalkblue}{RGB}{0, 105, 185}
\definecolor{chalkred}{RGB}{205, 35, 35}
\definecolor{chalkpink}{RGB}{230, 90, 120}
\definecolor{chalkcyan}{RGB}{15, 160, 190}
\definecolor{chalkgreen}{RGB}{20, 130, 50}
\definecolor{boxbg}{RGB}{252, 252, 254}
\definecolor{boxframe}{RGB}{150, 160, 180}

\pagestyle{fancy}
\fancyhf{}
\renewcommand{\headrulewidth}{0pt}
\renewcommand{\footrulewidth}{0pt}

\setlength{\parindent}{0pt}
\setlength{\parskip}{2.5pt}

\begin{document}

% ==============================================================================
% PAGE 1
% ==============================================================================
\noindent \textbf{\large Math 105L} \hfill \textbf{\large Lecture 9-1}\\[1pt]
{\textbf{\Large Differentiating Exponential Functions}}\\[-5pt]
\rule{\textwidth}{1.2pt}

\vspace{0.03in}

\textbf{\large Back to Calculus\dots}

\begin{itemize}[leftmargin=1.2em, topsep=1pt, itemsep=2pt]
  \item Given a function $f(x)$, its \textbf{derivative function} is given by the limit of the difference quotient:
  \[
    f'(x) = \lim_{h \to 0} \frac{f(x+h) - f(x)}{h}
  \]
  \item To approximate the derivative of a function at a point $x = a$, we plug in a small non-zero value of $h$ into the difference quotient:
  \[
    f'(a) \approx \frac{f(a+h) - f(a)}{h}
  \]
  \item \textbf{Proportionality Definition:} The quantity $A$ is \textit{directly proportional} to the quantity $B$ if $A = k \cdot B$ for some constant $k$.
\end{itemize}

\vspace{0.03in}

\textbf{\large Preliminaries: Playing with Derivatives via Numerical Exploration}

Let $f(x) = 2^x$ and $g(x) = 3^x$. Using a small step $h = 0.01$, we compute the ratio of the derivative to the function value:
\begin{itemize}[leftmargin=1.2em, label=$\bullet$, topsep=1pt, itemsep=2pt]
  \item \textbf{Sample calculation at $x = -4$ for $f(x) = 2^x$:}
  \[
    f(-4) = 2^{-4} = 0.0625, \quad f'(-4) \approx \frac{2^{-4 + 0.01} - 2^{-4}}{0.01} \approx 0.0435 \implies \frac{f'(-4)}{f(-4)} \approx \frac{0.0435}{0.0625} \approx 0.696
  \]
  \item Across all $x$, the ratio $\frac{f'(x)}{f(x)}$ is constantly $\approx 0.693$, which yields:
  \[
    \frac{\frac{d}{dx} 2^x}{2^x} \approx 0.693 \implies \mathbf{\frac{d}{dx} 2^x \approx 0.693 \cdot 2^x}
  \]
  \item Similarly for $g(x) = 3^x$, the ratio $\frac{g'(x)}{g(x)}$ is constantly $\approx 1.0987$ (or $1.105$ for $h=0.01$):
  \[
    \mathbf{\frac{d}{dx} 3^x \approx 1.0987 \cdot 3^x}
  \]
  \item \textbf{Key Observation:} In both cases, the derivative of an exponential function is \textbf{directly proportional to the function itself!}
\end{itemize}

\vspace{0.03in}

\textbf{\large Examining the Evidence: Why Does This Happen?}

Two fundamental questions arise:
\begin{enumerate}[leftmargin=1.4em, topsep=1pt, itemsep=1pt]
  \item \textit{Why is the derivative of an exponential function directly proportional to the function itself?}
  \item \textit{What is the exact constant of proportionality?}
\end{enumerate}

Let's compute $\frac{d}{dx} 2^x$ directly using the limit definition of the derivative:
\begin{align*}
  \frac{d}{dx} 2^x &= \lim_{h \to 0} \frac{2^{x+h} - 2^x}{h} = \lim_{h \to 0} \frac{2^x \cdot 2^h - 2^x}{h} \\
  &= \lim_{h \to 0} \frac{2^x(2^h - 1)}{h} = 2^x \cdot \lim_{h \to 0} \frac{2^h - 1}{h}
\end{align*}

Because $2^x$ has no $h$ dependency, it pulls completely out of the limit!
\begin{itemize}[leftmargin=1.2em, label=$\bullet$, topsep=1pt, itemsep=1pt]
  \item Estimating the remaining limit by plugging in a small value $h = 0.001$:
  \[
    \lim_{h \to 0} \frac{2^h - 1}{h} \approx \frac{2^{0.001} - 1}{0.001} \approx 0.69315
  \]
  \item \textbf{General Base $a$:} Repeating for any base $a > 0$:
  \[
    \frac{d}{dx} a^x = \lim_{h \to 0} \frac{a^{x+h} - a^x}{h} = a^x \cdot \lim_{h \to 0} \frac{a^h - 1}{h}
  \]
\end{itemize}

\newpage

% ==============================================================================
% PAGE 2
% ==============================================================================
\noindent \textbf{\large Math 105L} \hfill \textbf{\large Lecture 9-1 (Page 2)}\\[1pt]
{\textbf{\Large A Pretty (and Important!) Derivative \& The Natural Base $e$}}\\[-5pt]
\rule{\textwidth}{1.2pt}

\vspace{0.03in}

\begin{tcolorbox}[colback=boxbg, colframe=boxframe, boxrule=0.8pt, arc=3pt, left=8pt, right=8pt, top=4pt, bottom=4pt]
\textbf{General Proportionality Law for Exponential Derivatives:}\\[2pt]
\[
  \frac{d}{dx} a^x = k \cdot a^x \quad \text{where} \quad k = \lim_{h \to 0} \frac{a^h - 1}{h}
\]
\end{tcolorbox}

\vspace{0.02in}

\textbf{\large Finding the Special Base Where $k = 1$}

Suppose we could find a special base $a$ such that the constant of proportionality is exactly $1$ ($k = 1$). For that magic base $a$, we would have:
\[
  \mathbf{\frac{d}{dx} a^x = a^x}
\]
\begin{itemize}[leftmargin=1.2em, label=$\bullet$, topsep=1pt, itemsep=2pt]
  \item When $a = 2 \implies k \approx 0.693 < 1$.
  \item When $a = 3 \implies k \approx 1.099 > 1$.
  \item By the Intermediate Value Property, such an $a$ must lie between $2$ and $3$:
  \begin{itemize}[label=--]
    \item Try $a = 2.5 \implies \frac{2.5^{0.001} - 1}{0.001} \approx 0.9166$.
    \item Try $a = 2.718 \implies \frac{2.718^{0.001} - 1}{0.001} \approx 0.99989 \approx 1$.
  \end{itemize}
  \item \textbf{Analytical Derivation of the Base $e$:}
  \[
    \frac{a^h - 1}{h} \approx 1 \implies a^h - 1 \approx h \implies a^h \approx 1 + h \implies a \approx (1 + h)^{1/h}
  \]
  Taking the limit as $h \to 0$, and setting $n = \frac{1}{h}$ (so $h = \frac{1}{n}$ as $n \to \infty$):
  \[
    a = \lim_{h \to 0} (1 + h)^{1/h} = \lim_{n \to \infty} \left(1 + \frac{1}{n}\right)^n = \mathbf{e} \approx 2.7182818\dots
  \]
\end{itemize}

\begin{tcolorbox}[colback=chalkblue!6, colframe=chalkblue, boxrule=1pt, arc=3pt, left=8pt, right=8pt, top=5pt, bottom=5pt]
\textbf{\large Derivative of $e^x$:}
\[
  \mathbf{\frac{d}{dx} e^x = e^x}
\]
The natural exponential function is its own derivative! At every point $(x, e^x)$, the slope of the tangent line equals the $y$-value.
\end{tcolorbox}

\vspace{0.04in}

\textbf{\large Finding the General Formula: Derivative of $a^x$ via Chain Rule}

\begin{itemize}[leftmargin=1.2em, topsep=1pt, itemsep=2pt]
  \item \textbf{Preliminary Chain Rule Question:} What is $\frac{d}{dx} e^{kx}$?
  \[
    \text{inner: } u = kx \implies u' = k, \quad \text{outer: } e^u \implies (e^u)' = e^u
  \]
  \[
    \frac{d}{dx} e^{kx} = e^{kx} \cdot k = \mathbf{k e^{kx}}
  \]
  \item \textbf{Rewriting Base $a$ into Base $e$:}
  Recall from algebra that $a = e^{\ln a}$. Therefore:
  \[
    a^x = (e^{\ln a})^x = e^{(\ln a)x} \quad (\text{here } k = \ln a)
  \]
  \item \textbf{Applying the Chain Rule:}
  \[
    \frac{d}{dx} a^x = \frac{d}{dx} e^{(\ln a)x} = (\ln a) \cdot e^{(\ln a)x} = \mathbf{(\ln a) a^x}
  \]
\end{itemize}

\begin{tcolorbox}[colback=chalkgreen!6, colframe=chalkgreen, boxrule=1pt, arc=3pt, left=8pt, right=8pt, top=5pt, bottom=5pt]
\textbf{\large General Exponential Derivative Rule:}
\[
  \mathbf{\frac{d}{dx} a^x = (\ln a) \cdot a^x}
\]
\textbf{Examples from Class:}
\[
  \frac{d}{dx} 2^x = (\ln 2) \cdot 2^x \approx 0.69315 \cdot 2^x, \qquad \frac{d}{dx} 3^x = (\ln 3) \cdot 3^x \approx 1.0986 \cdot 3^x
\]
\end{tcolorbox}

\end{document}

